\section{Gravity bridge: contraction channel from transverse excitation}
\label{sec:gravity}

This bridge imports the near-identity metric expansion \eqref{eq:metric-near-identity} and
the quasi-static in-brane equilibrium from Paper~I (\Cref{subsec:gravity-channel}) and asks
whether the resulting slowly varying contraction field can serve as the quantitative substrate
of a Newtonian gravitational potential.
The geometric mechanism is fixed by Paper~I; the open problem addressed here is its quantitative
reduction.

\subsection{From the imported contraction channel to an explicit forcing problem}

Paper~I (\Cref{subsec:gravity-channel}) establishes the near-identity metric expansion
\eqref{eq:metric-near-identity}, the small-slope in-brane strain \eqref{eq:strain-near-identity},
and that the transverse term $\partial_i u\,\partial_j u$ acts as a source for the slow in-brane
field $\vec\xi$ through quasi-static balance $\partial_j\sigma_{ij}=0$ (the gravity channel
requires the in-brane relaxation $\vec\xi\neq 0$).
Writing that balance out with the constitutive law
$\sigma_{ij}=\lambda\,\mathrm{tr}(\varepsilon)\delta_{ij}+2\mu\,\varepsilon_{ij}$ gives the
explicit forcing/eigenstrain problem for $\vec\xi$:
\begin{equation}
  \partial_j\bigl[\lambda\,(\partial_k\xi_k)\delta_{ij}+\mu(\partial_j\xi_i+\partial_i\xi_j)\bigr]
  = -\tfrac12\partial_j[\lambda\,|\nabla u|^2\delta_{ij}+\mu(\partial_i u\partial_j u+\partial_j u\partial_i u)].
  \label{eq:gravity-forcing}
\end{equation}
The right-hand side is the effective body force sourced by a transverse profile $u(x)$.
The solution $\vec\xi$ is the in-brane contraction field; its divergence $\partial_i\xi_i$ is
the local volume contraction (dilation) pattern.

The sign of this response is fixed by Paper~I (\Cref{subsec:gravity-channel}): under pre-tension
($\alpha<1$) the added transverse arc length raises tensile stress, so the surrounding brane is
pulled \emph{toward} the excitation (contraction) rather than pushed outward, with no separate
assumption required.

\subsection{Unification: length contraction and time dilation}

Because the underlying Pythagorean mechanism is direction-blind, Paper~I
(\Cref{subsec:gravity-channel}) shows that the same rule applied symmetrically across the four
ambient directions makes spatial contraction (timelike-direction excitation lengthening spacelike
links --- the channel quantified in this section) and gravitational time dilation
(spacelike-direction excitation lengthening the timelike link) two faces of one rule.
The present section quantifies only the first face.

\subsection{Contraction field as a gravity-like background}

Linear wave packets propagating on top of the slowly varying $\vec\xi$-background experience
branch-dependent effective coefficients in the linearized wave operator.
In the effective-metric viewpoint (Bridge~1, \Cref{sec:lorentz}), the $\vec\xi$-sourced
dilation pattern acts as an analogue curved spacetime that focuses or defocuses wave transport.

Concretely, a transverse excitation concentrated in a region $\mathcal{V}$ produces through
\eqref{eq:gravity-forcing} an outward-decaying contraction field $\partial_i\xi_i(x)$ that
attracts nearby wave packets inward.
In the weak-field, slowly varying limit, this biases propagation in precisely the way a
Newtonian potential would.

\paragraph{Weak equivalence principle as a corollary.}
Paper~IV identifies the inertial mass of a localized mode with its rest energy,
$m_{\mathrm{inertial}}=E_{\mathrm{rest}}/c^2$ \citep{Einstein1905inertia}.
The gravity channel developed here sources the effective potential $\Phi$ from the
\emph{same} localized excitation energy: the body force on the right-hand side of
\eqref{eq:gravity-forcing} is set by the transverse energy density $\tfrac{1}{2}|\nabla u|^2$.
The gravitational charge density $\rho_{\mathrm{excitation}}$ in \eqref{eq:gravity-poisson-target}
and the inertial mass density are therefore the same substrate quantity, so the equality of
gravitational and inertial mass --- the weak equivalence principle --- would hold by construction
rather than as an independent postulate.
This is a structural prediction contingent on two results not yet closed: the inertial-mass
identification of Paper~IV (gated on its explicit toroidal energy integral and the classical
$4/3$ consistency test) and verification, in the Newtonian reduction below, that the Poisson
source $\rho_{\mathrm{excitation}}$ is precisely the total rest-energy density of the mode.
It is directly testable: in the probe-packet deflection test (\Cref{subsec:gravity-open}),
all modes must fall identically in a given contraction background, independent of internal
composition.

\subsection{Open quantitative reduction: Newtonian limit}
\label{subsec:gravity-open-reduction}

What remains is the quantitative reduction.
The goal is to derive a scalar potential $\Phi(x)$ from the coarse-grained in-brane strain or
elastic energy density such that
\begin{equation}
  \nabla^2 \Phi = 4\pi G_{\text{eff}}\,\rho_{\text{excitation}},
  \label{eq:gravity-poisson-target}
\end{equation}
with a coupling $G_{\text{eff}}$ expressed in substrate parameters $(k_s, a, \alpha, m)$.
The specific steps required:

\begin{enumerate}
  \item \textbf{Solve the forcing problem \eqref{eq:gravity-forcing} in a scalar approximation.}
  For a spherically symmetric transverse profile $u(r)$, symmetry reduces
  \eqref{eq:gravity-forcing} to a scalar ODE for $\partial_i\xi_i$.

  \item \textbf{Identify the effective potential.}
  Show that $\Phi \propto \partial_i\xi_i$ (or a suitable Green's function convolution thereof)
  satisfies a Poisson-type equation with a source proportional to the transverse energy density
  $\tfrac12|\nabla u|^2$.

  \item \textbf{Match to Newtonian gravity.}
  Express $G_{\text{eff}}$ in terms of $k_s$, $a$, $\alpha$, $m$; derive the dimensional
  prediction $G_{\text{eff}} \sim k_s a^5 / m^2$ (or similar) that sets the scale of the
  coupling.

  \item \textbf{Numerical check.}
  Simulate a localized transverse excitation, measure the induced $\vec\xi$-field, and compare
  its spatial decay to the Green's function prediction.
  Inject a long-wavelength probe packet and measure whether its trajectory is deflected inward
  at the rate predicted by $\Phi$.
\end{enumerate}

\subsection{Gravity bridge priority open problems}
\label{subsec:gravity-open}

\begin{enumerate}
  \item \textbf{Scalar reduction of \eqref{eq:gravity-forcing}.}
  Derive the spherically symmetric Poisson equation for the contraction field and obtain
  $G_{\text{eff}}$ in closed form.

  \item \textbf{Contraction-channel extraction (numerical).}
  Create a localized transverse Gaussian excitation; measure the induced slow
  $\partial_i\xi_i$-field; compare its spatial profile to the prediction from step 1.

  \item \textbf{Probe-packet deflection test.}
  Inject a long-wavelength wave packet, measure the lateral deflection, and check that the
  deflection angle matches the effective potential derived in step 1; this also serves as a
  universality-of-free-fall check --- modes of different internal composition must fall
  identically in a given contraction background, as required by the weak equivalence principle
  corollary above.

  \item \textbf{4D symmetric treatment.}
  Extend the slice-restricted analysis to a 4D treatment of the brane action, expressing both
  length contraction and time dilation in a single expression, and derive the full stress-energy
  coupling tensor.
\end{enumerate}